% Options for packages loaded elsewhere
\PassOptionsToPackage{unicode}{hyperref}
\PassOptionsToPackage{hyphens}{url}
\PassOptionsToPackage{dvipsnames,svgnames,x11names}{xcolor}
%
\documentclass[
  13pt,
  a4paperpaper,
  12pt]{article}
\usepackage{amsmath,amssymb}
\usepackage{setspace}
\usepackage{iftex}
\ifPDFTeX
  \usepackage[T1]{fontenc}
  \usepackage[utf8]{inputenc}
  \usepackage{textcomp} % provide euro and other symbols
\else % if luatex or xetex
  \usepackage{unicode-math} % this also loads fontspec
  \defaultfontfeatures{Scale=MatchLowercase}
  \defaultfontfeatures[\rmfamily]{Ligatures=TeX,Scale=1}
\fi
\usepackage{lmodern}
\ifPDFTeX\else
  % xetex/luatex font selection
  \setmainfont[]{Times New Roman}
  \setmonofont[]{Courier New}
\fi
% Use upquote if available, for straight quotes in verbatim environments
\IfFileExists{upquote.sty}{\usepackage{upquote}}{}
\IfFileExists{microtype.sty}{% use microtype if available
  \usepackage[]{microtype}
  \UseMicrotypeSet[protrusion]{basicmath} % disable protrusion for tt fonts
}{}
\makeatletter
\@ifundefined{KOMAClassName}{% if non-KOMA class
  \IfFileExists{parskip.sty}{%
    \usepackage{parskip}
  }{% else
    \setlength{\parindent}{0pt}
    \setlength{\parskip}{6pt plus 2pt minus 1pt}}
}{% if KOMA class
  \KOMAoptions{parskip=half}}
\makeatother
\usepackage{xcolor}
\usepackage[margin=1.5cm,top=1.5cm,bottom=1.5cm,left=1.5cm,right=1.5cm,includeheadfoot]{geometry}
\usepackage{listings}
\newcommand{\passthrough}[1]{#1}
\lstset{defaultdialect=[5.3]Lua}
\lstset{defaultdialect=[x86masm]Assembler}
\setlength{\emergencystretch}{3em} % prevent overfull lines
\providecommand{\tightlist}{%
  \setlength{\itemsep}{0pt}\setlength{\parskip}{0pt}}
\setcounter{secnumdepth}{3}
% LaTeX Preamble for Insect Perception Research
\documentclass[13pt,a4paper]{article}
\usepackage[utf8]{inputenc}
\usepackage[T1]{fontenc}
\usepackage{geometry}
\usepackage{graphicx}
\usepackage{amsmath}
\usepackage{amsfonts}
\usepackage{amssymb}
\usepackage{booktabs}
% \usepackage{hyperref}  % Commented out to prevent unwanted metadata
\usepackage{color}
\usepackage{xcolor}
\usepackage{listings}
\usepackage{fancyvrb}
\usepackage{setspace}

% Page geometry - more accessible margins
\geometry{margin=1in, top=1in, bottom=1in}

% Hyperref package removed to prevent unwanted metadata

% Custom colors
\definecolor{codebg}{RGB}{245, 245, 245}
\definecolor{codeborder}{RGB}{200, 200, 200}
\definecolor{codefg}{RGB}{50, 50, 50}

% Code listing setup
\lstset{
    backgroundcolor=\color{codebg},
    basicstyle=\normalsize\ttfamily\color{codefg},
    breakatwhitespace=false,
    breaklines=true,
    captionpos=b,
    frame=single,
    framerule=0.5pt,
    framesep=5pt,
    rulecolor=\color{codeborder},
    numbers=left,
    numbersep=5pt,
    numberstyle=\small\color{codefg}
}

% Font and spacing for accessibility
% Using standard LaTeX fonts for better compatibility
\renewcommand{\familydefault}{\rmdefault}
\renewcommand{\sfdefault}{\rmdefault}

% Single spacing for academic format
\setstretch{1.0}
\setlength{\parindent}{0.5in}
\setlength{\parskip}{0em}

% Ensure left-aligned text (not centered)
\raggedright

% Override any centering for academic documents
\makeatletter
\renewcommand{\maketitle}{%
  \begin{titlepage}%
    \let\footnotesize\small
    \let\footnoterule\relax
    \let\footnote\thanks
    \begin{center}%
      {\LARGE \@title \par}%
      \vskip 1.5em%
      {\large
        \lineskip .5em%
        \begin{tabular}[t]{c}%
          \@author
        \end{tabular}\par}%
      \vskip 1em%
      {\large \@date}%
    \end{center}%
    \par
    \@thanks
  \end{titlepage}%
  \setcounter{footnote}{0}%
  \global\let\thanks\relax
  \global\let\maketitle\relax
  \global\let\@thanks\@empty
  \global\let\@author\@empty
  \global\let\@title\@empty
  \global\let\@date\@empty
  \global\let\title\relax
  \global\let\author\relax
  \global\let\date\relax
  \global\let\and\relax
}
\makeatother

% Ensure body text is left-aligned
\raggedright
\setlength{\parindent}{0.5in}

% Comprehensive numbering for figures, tables, and equations
\usepackage{chngcntr}
\counterwithout{figure}{section}
\counterwithout{table}{section}
\counterwithout{equation}{section}

% Figure and table caption formatting
\usepackage[font=small,labelfont=bf,textfont=it]{caption}
\captionsetup[figure]{position=bottom,skip=10pt}
\captionsetup[table]{position=top,skip=10pt}

% Equation numbering and formatting
\usepackage{amsthm}
\newtheorem{theorem}{Theorem}[section]
\newtheorem{lemma}[theorem]{Lemma}
\newtheorem{proposition}[theorem]{Proposition}
\newtheorem{corollary}[theorem]{Corollary}
\newtheorem{definition}[theorem]{Definition}
\newtheorem{example}[theorem]{Example}
\newtheorem{remark}[theorem]{Remark}

% Section formatting - ensure proper academic style
\usepackage{titlesec}
\titlespacing{\section}{0pt}{12pt}{6pt}
\titlespacing{\subsection}{0pt}{10pt}{4pt}
\titlespacing{\subsubsection}{0pt}{8pt}{3pt}

% Ensure sections are left-aligned and properly formatted
\titleformat{\section}{\large\bfseries}{\thesection}{1em}{}
\titleformat{\subsection}{\normalsize\bfseries}{\thesubsection}{1em}{}
\titleformat{\subsubsection}{\normalsize\itshape}{\thesubsubsection}{1em}{}

% Title and author
\title{When do bugs see (infra)red? On the Visual and Infra-red in the Insect Perceptual Apparatus}
\author{Tucker Chambers \\
\small Email: tucker.chambers@example.com \\
\small ORCID: 0000-0000-0000-0000 \and 0000-0000-0000-0001 \\
Daniel A. Friedman \\
\small Email: daniel.friedman@example.com}
\ifLuaTeX
  \usepackage{selnolig}  % disable illegal ligatures
\fi
\IfFileExists{bookmark.sty}{\usepackage{bookmark}}{\usepackage{hyperref}}
\IfFileExists{xurl.sty}{\usepackage{xurl}}{} % add URL line breaks if available
\urlstyle{same}
\hypersetup{
  colorlinks=true,
  linkcolor={blue},
  filecolor={blue},
  citecolor={blue},
  urlcolor={blue},
  pdfcreator={LaTeX via pandoc}}

\author{}
\date{}

\begin{document}

{
\hypersetup{linkcolor=black}
\setcounter{tocdepth}{3}
\tableofcontents
}
\setstretch{1.0}
\section{Introduction}\label{sec:introduction}

Olfaction, the ability to detect and identify airborne molecules,
represents one of the most fundamental sensory modalities across
biological systems. This sense exhibits remarkable conservation across
diverse taxa, with insects demonstrating particularly sophisticated
chemosensory capabilities that challenge traditional mechanistic
explanations.

\subsection{Current Understanding of Insect
Olfaction}\label{current-understanding-of-insect-olfaction}

The classical stereochemical theory of olfaction posits that molecular
recognition occurs through shape complementarity between odor molecules
and olfactory receptors (ORs) on the cellular membranes of olfactory
receptor neurons (ORNs). This lock-and-key mechanism initiates ionic
cascades that generate measurable neural responses within milliseconds
of stimulus detection.

\textbf{Receptor Diversity and Specificity}: Insects possess hundreds of
distinct OR types, yet can discriminate among billions of perceptible
odors. This remarkable capability is achieved through combinatorial
activation patterns, where individual odors activate multiple receptors
with varying affinities, creating high-dimensional neural
representations despite individual receptor broad-tuning.

\subsection{Limitations of Stereochemical
Theory}\label{limitations-of-stereochemical-theory}

\subsubsection{Isotope Discrimination
Evidence}\label{isotope-discrimination-evidence}

The stereochemical theory faces significant challenges from isotope
discrimination studies. Molecules with identical shapes and chemical
structures but different isotopic compositions elicit distinct olfactory
responses, demonstrating that molecular geometry alone cannot explain
odor discrimination.

\textbf{Quantitative Evidence}: Studies on \emph{Drosophila
melanogaster} show that deuterated homologues of known odorants produce
unique behavioral responses despite maintaining identical molecular
shapes. This finding is quantitatively supported by vibrational
spectroscopy, where deuteration shifts infrared emission spectra by 2-3
μm while preserving molecular geometry.

\subsubsection{Response Time
Inconsistencies}\label{response-time-inconsistencies}

Traditional molecular binding models cannot account for the extremely
rapid response times observed in insect olfaction. Insect ORNs
demonstrate response latencies of 1-5 ms, comparable to photoreceptor
(0.1 ms) and auditory receptor (0.16 ms) response times.

\textbf{Mechanistic Implications}: These rapid responses suggest
detection mechanisms that bypass the slower processes of molecular
diffusion, pore penetration, and receptor binding. The observed response
times are more consistent with electromagnetic wave detection than with
molecular binding kinetics.

\subsection{The Vibrational Theory
Alternative}\label{the-vibrational-theory-alternative}

The vibrational theory of olfaction proposes that insects detect the
unique electromagnetic radiation emitted by free-floating odor molecules
rather than relying solely on geometric or chemical information at
receptor binding surfaces.

\subsubsection{Atmospheric Transmission
Windows}\label{atmospheric-transmission-windows}

A compelling aspect of the vibrational theory is the correspondence
between atmospheric transmission characteristics and semiochemical
emission spectra. Earth's atmosphere exhibits specific transmission
windows in the mid- and long-infrared ranges (2-5 μm, 8-14 μm, 17-25 μm)
that precisely overlap with the emission spectra of insect
semiochemicals.

\textbf{Environmental Optimization}: These transmission windows enable
long-range infrared detection, providing insects with access to
environmental information that would be unavailable through molecular
diffusion alone. The atmospheric transmission function is quantitatively
modeled in Section \ref{sec:mathematical_appendix}.

\subsubsection{Sensilla as Electromagnetic
Antennas}\label{sensilla-as-electromagnetic-antennas}

Insect antennae and sensilla exhibit morphological adaptations that
suggest optimization for electromagnetic detection. Sensilla dimensions
(6-160 μm for trichodea, 2-8 μm for basiconica) correspond closely to
the wavelengths of infrared radiation emitted by semiochemicals.

\textbf{Structural Evidence}: The porous architecture of sensilla,
traditionally interpreted as molecular transport conduits, may serve
dual functions as electromagnetic waveguides and molecular filters. This
interpretation is supported by dielectric waveguide theory and resonant
cavity modeling.

\subsection{Research Objectives and
Approach}\label{research-objectives-and-approach}

This paper examines the vibrational theory through multiple analytical
domains:

\begin{enumerate}
\def\labelenumi{\arabic{enumi}.}
\tightlist
\item
  \textbf{Morphological Analysis}: Quantitative assessment of sensilla
  dimensions and their correspondence to infrared wavelengths
\item
  \textbf{Neurological Investigation}: Analysis of response time data
  and neural encoding mechanisms
\item
  \textbf{Behavioral Studies}: Examination of insect responses to
  infrared stimuli and environmental conditions
\item
  \textbf{Spectroscopic Validation}: Measurement and analysis of
  semiochemical emission spectra
\item
  \textbf{Computational Modeling}: Implementation of theoretical
  frameworks in tested computational models
\end{enumerate}

\subsection{Theoretical Framework
Integration}\label{theoretical-framework-integration}

The vibrational theory integrates multiple physical principles:

\begin{itemize}
\tightlist
\item
  \textbf{Electromagnetic Wave Theory}: Maxwell's equations and
  waveguide propagation in dielectric media
\item
  \textbf{Quantum Mechanics}: Electron tunneling and phonon coupling in
  olfactory receptors
\item
  \textbf{Resonant Cavity Theory}: Sensilla as frequency-tuned
  electromagnetic resonators
\item
  \textbf{Piezoelectric Effects}: Microtubule arrays as signal
  amplifiers and frequency selectors
\end{itemize}

\subsection{Empirical Validation
Strategy}\label{empirical-validation-strategy}

All theoretical predictions are implemented in tested computational
models that generate quantitative predictions for experimental
validation. The mathematical framework presented in Section
\ref{sec:mathematical_appendix} provides specific equations that can be
tested through:

\begin{itemize}
\tightlist
\item
  \textbf{Sensilla Response Measurements}: Direct testing of infrared
  sensitivity across different frequencies
\item
  \textbf{Behavioral Assays}: Quantification of insect responses to
  infrared stimuli
\item
  \textbf{Neural Recording}: Measurement of ORN responses to
  electromagnetic stimulation
\item
  \textbf{Environmental Studies}: Analysis of atmospheric transmission
  effects on detection range
\end{itemize}

This integrated approach ensures that theoretical predictions are
grounded in empirical reality and provides a framework for future
experimental validation of the vibrational theory.

\end{document}
